% Hand-written narrative; numbers are macros.
All detectors except the signature rules find the three S3 administrators who log in as their
supervisor (\DROursSThree{} for ours): new-PC logons, after-hours activity and a burst of e-mail make
those days stand out. About half of the S1 insiders, who copy data to USB after hours and visit a leak
site, are caught by every method (\DROursSOne{} for ours, \DRIFSOne{} for Isolation Forest). S2,
month-long job hunting followed by data theft, is where the methods separate: Isolation Forest
catches \DRIFSTwo{} of these insiders, BudgetAwareConformal \DROursSTwo{} and LOF \DRLOFSTwo. Their
individual days look ordinary to a population model; they stand out only against the user's own
history, and only on a few days of their month-long malicious period.

The contrast between S2 results and S3 results isolates the value of per-user baselines: when the
attack is extreme relative to any user's history (S3: logging in as a different account on a PC
never seen before), a population model works as well as a user-specific one. When the attack
is a small change in a user's normal pattern (S2: slightly more job-site visits, a USB connection
on an unusual day), only the per-user baseline has the resolution to detect it. This finding is
consistent with the ablation, where removing C1 (the per-user baseline) collapses the detection
rate to the Isolation Forest level even though all other components remain intact.
